%%%%%%%%%%%%%%%%%%%%%%%%%%%%%%%%%%%%%%%%%
% Focus Beamer Presentation
% LaTeX Template
% Version 1.0 (8/8/18)
%
% This template has been downloaded from:
% http://www.LaTeXTemplates.com
%
% Original author:
% Pasquale Africa (https://github.com/elauksap/focus-beamertheme) with modifications by 
% Vel (vel@LaTeXTemplates.com)
%
% Template license:
% GNU GPL v3.0 License
%
% Important note:
% The bibliography/references need to be compiled with bibtex.
%
%%%%%%%%%%%%%%%%%%%%%%%%%%%%%%%%%%%%%%%%%

%----------------------------------------------------------------------------------------
%	PACKAGES AND OTHER DOCUMENT CONFIGURATIONS
%----------------------------------------------------------------------------------------

\documentclass{beamer}
\setbeamertemplate{caption}[numbered]

\usetheme{focus} % Use the Focus theme supplied with the template
% Add option [numbering=none] to disable the footer progress bar
% Add option [numbering=fullbar] to show the footer progress bar as always full with a slide count

% Uncomment to enable the ice-blue theme
%\definecolor{main}{RGB}{92, 138, 168}
%\definecolor{background}{RGB}{240, 247, 255}

%------------------------------------------------

\usepackage{booktabs} % Required for better table rules

%----------------------------------------------------------------------------------------
%	 TITLE SLIDE
%----------------------------------------------------------------------------------------

\title{WiFinger: Talk to Your Smart Devices with Finger-grained Gesture}

\subtitle{Qualifying Examination}

\author{Rishiraj Adhikary}

\titlegraphic{\includegraphics[scale=0.20]{Images/iitgn.png}} % Optional title page image, comment this line to remove it

\institute{Computer Science and Engineering \\ IIT Gandhinagar}

\date{12 03 2021}

%------------------------------------------------

\begin{document}

%------------------------------------------------

\begin{frame}
	\maketitle % Automatically created using the information in the commands above
\end{frame}

\begin{frame}{Overview}
\tableofcontents
\end{frame}

%----------------------------------------------------------------------------------------
%	 SECTION 1
%----------------------------------------------------------------------------------------

\section{Introduction} % Section title slide, unnumbered

%\begin{frame}{Introduction}
%WiFi Signals Can Sense People's Location and Activities
%	\begin{center}
%		\begin{figure}
%			\includegraphics[width=1.0\textwidth]{Images/wikey.png}
%			\caption{Detecting keyboard strokes using WiFi~\cite{ali2015keystroke}}
%		\end{figure}
%     \end{center}
%\end{frame}

\begin{frame}{Introduction}
WiFi Signals Can Sense People's Location and Activities
	\begin{center}
		\begin{figure}
			\includegraphics[width=0.5\textwidth]{Images/wiaction.png}
			\caption{Detecting gestures using WiFi~\cite{pu2013whole}}
		\end{figure}
     \end{center}
\end{frame}

\subsection{Objective}

\begin{frame}{Objective}
	\begin{columns}
		\begin{column}{0.5\textwidth}
  			\emph{WiFinger is a wireless system that utilizes commercial WiFi devices to achieve human-computer interaction by recognizing people's finger-grained gestures.} 
		\end{column}
		\begin{column}{0.5\textwidth}  %%<--- here
    		\begin{center}
    			\begin{figure}
     				\includegraphics[width=0.75\textwidth]{Images/wifinger.jpg}
     				\caption{Demonstration of WiFinger~\cite{li2016wifinger}}
     			\end{figure}
     		\end{center}
		\end{column}
	\end{columns}
\end{frame}

\begin{frame}{Objective}
	\begin{center}
		\begin{figure}
			\includegraphics[width=0.75\textwidth]{Images/asl.png}
			\caption{Parts of finger gestures that WiFinger can detect and recognize~\cite{li2016wifinger}}
		\end{figure}
     \end{center}
\end{frame}

\subsection{Related Work}
\begin{frame}{Related Work}
Prior (and current) work on gesture detection can be categorized into two groups.
\begin{itemize}
	\item Device based
		\begin{itemize}
			\item Audio based
			\item Vision based
			\item Sensor based
		\end{itemize}
	\item Device free - Wireless Signal
\end{itemize}

\end{frame}

\subsection{Background}
\begin{frame}{Background}
\textbf{Received Signal Strength (RSS)}:\\
Universal Software Radio Peripheral captures RSS values from WiFi signals~\cite{sigg2013leveraging, sigg2013rf}. RSS values are not suitable for recognizing fine-grained motions such as gestures in standard American Sign Language (ASL).
\end{frame}

\begin{frame}{Background}
\textbf{Channel State Information (CSI)}\\
	\begin{itemize}
		\item CSI refers to known channel properties of a communication link. The channel between transmitter and receiver comprises of multiple subcarriers. 
		\item $y$ is the received vector and $x$ is the transmitted vector,
		\begin{equation*}
			y=Hx + n
		\end{equation*}
		\item $n$ is the noise vector. $H$ is the channel frequency response. 
	\end{itemize}
\end{frame}

\begin{frame}{Background}
\textbf{Channel State Information (CSI)}\\	
	\begin{equation*}
		y=Hx + n
	\end{equation*}
The dimension of $H$ is $N_c$ x $N_t$ x $N_r$.\\
\vspace{10px}
$N_c$: Number of sub carriers. $N_t$: Number of transmit antennas. $N_r$: Number of receive antennas.\\
\begin{figure}[H]
		\centering
		\includegraphics[width=0.4\textwidth]{Images/csi.png}
		\caption{Subcarrier-level signal strength computed from channel state information for four single-antenna 802.11n links~\cite{halperin2011tool}.}
	\end{figure}
\end{frame}




%------------------------------------------------




\section{Approach}

\subsection{Overview}
\begin{frame}{Overview}
	\begin{center}
		\begin{figure}
			\includegraphics[width=1.\textwidth]{Images/overview.png}
			\caption{Framework of WiFinger~\cite{li2016wifinger}}
		\end{figure}
     \end{center}
\end{frame}


\subsection{Signal Preprocessing}
\begin{frame}{Signal Preprocessing}
Signal changes caused by finger motions lie at the low end of the frequency spectrum while noise induced by hardware imperfections has a relatively high frequency.
	\begin{center}
		\begin{figure}
			\includegraphics[width=1.\textwidth]{Images/quals-part-1-denosing.png}
			\caption{Noise reduction using low pass filter~\cite{li2016wifinger}}
		\end{figure}
     \end{center}
\end{frame}

\subsection{Gesture Extraction}
\begin{frame}{Gesture Detection}
\textbf{Preprocessing}
	\begin{columns}
		\column{0.5\textwidth}
			\begin{itemize}
				\item The CSI stream is cut into bins using a sliding window
				\item The window size is 500. 
				\item Each bin is a matrix of size 30 x 500 $ = \mathbf{M}$
			\end{itemize}
		
		\column{0.5\textwidth}
			\begin{figure}
				\includegraphics[width=1\linewidth]{Images/quals-part-1-gesture-d.png}
			\caption{Windowing of the CSI~\cite{li2016wifinger}}
			\end{figure}	
	\end{columns}
\end{frame}

\begin{frame}{Gesture Detection}
\textbf{Correlation Estimation}
	\begin{columns}
		\column{0.5\textwidth}
			\begin{itemize}
				\item WiFinger calculates the correaltion matrix as $\mathbf{M^{T}}$ x $\mathbf{M}$
				\item The value of the second eigenvector of the above matrix indicates the presence and absence of a \emph{sign}
			\end{itemize}
		\column{0.5\textwidth}
			\begin{figure}
			\includegraphics[width=1\linewidth]{Images/quals-part-1-gesture-d-2.png}
			\caption{Gesture detected in 2 windows~\cite{li2016wifinger}}
			\end{figure}
	\end{columns}

\end{frame}

\subsection{Feature Extraction}
\begin{frame}{Feature Extraction}
	\begin{center}
		\begin{figure}
			\includegraphics[width=0.75\textwidth]{Images/quals-part-2-gesture-d-3.png}
			\caption{Gesture Profile Extraction~\cite{li2016wifinger}}
		\end{figure}
    \end{center}
The profile of a particular sign can be mathematically represented as $\mathbf{P_i}$ = $[H_{t_{i}^{s}} \ldots H_{t_{i}^{e}}]$
\end{frame}

\begin{frame}{Feature Extraction}
	\begin{center}
		\begin{figure}
			\includegraphics[width=1.\textwidth]{Images/quals-part-2-subcarriers.png}
			\caption{Gesture Feature Extraction~\cite{li2016wifinger}}
		\end{figure}
    \end{center}
WiFinger combines 30 subcarriers by averaging every 6 subcarriers and then concatenates them to form a synthetic waveform.
\end{frame}

\begin{frame}{Feature Extraction}
WiFinger compresses the feature vectors by utilizing \textbf{Discrete Wavelet Transform} (DWT).\\

	\begin{itemize}
		\item Reduces computational cost compared to Fast Fourier Transform (FFT).
		\item Preserves both time and frequency domain information.
	\end{itemize}
\end{frame}

\subsection{Classification}
\begin{frame}{Classification}
	\begin{itemize}
		\item WiFinger utilizes k-Nearest Neighbor (kNN) classifier to recognize different finger gestures. 
		\item Feature vector of gestures might not share the same length.
	\end{itemize}
\end{frame}

\begin{frame}{Dynamic Time Warping}
	\textbf{Dynamic Time Warping} (DTW) provides intuitive distance between two waveforms and can 	be resilient to signal distortion and shift.
	\begin{figure}
		\includegraphics[width=0.75\linewidth]{Images/dtw-1.png}
		\caption{kNN and kNN with DTW\footnote{https://github.com/markdregan/K-Nearest-Neighbors-with-Dynamic-Time-Warping}}
	\end{figure}
\end{frame}



%------------------------------------------------


\section{Evaluation and Result}
\begin{frame}{Implementation}
	\begin{center}
		\begin{figure}
			\includegraphics[width=1.\textwidth]{Images/impl.png}
			\caption{The experimental setup of WiFinger~\cite{li2016wifinger}}
		\end{figure}
    \end{center}
\end{frame}

\begin{frame}{Evaluation}
	\begin{itemize}
		\item 10 users volunteered for the study.
		\item 9 users performed each gestures 35 times. One user performed each gesture 70 times. 
	\end{itemize}
\end{frame}

\begin{frame}{Result}
	\begin{figure}
		\includegraphics[width=0.50\linewidth]{Images/gesture-detection-accuracy.png}
		\caption{Finger gesture extraction accuracy per gesture for users 1-10~\cite{li2016wifinger}}
	\end{figure}
	
		\begin{columns}
		\column{0.5\textwidth}
			\begin{figure}
				\includegraphics[width=0.75\linewidth]{Images/recognition-accuracy.png}
			\caption{Average recognition accuracy per gesture~\cite{li2016wifinger}}
			\end{figure}
		\column{0.5\textwidth}
			\begin{figure}
				\includegraphics[width=0.75\linewidth]{Images/recognition-accuracy-per-user.png}
			\caption{Average finger gesture recognition accuracy per user~\cite{li2016wifinger}}
			\end{figure}
	\end{columns}
\end{frame}

\section{Limitations and Summary}
\begin{frame}{Limitations}
	\begin{itemize}
		\item Requires line of sight between transmitter and receiver.
		\item Presence of human body motion, moving objects and the orientation of transceiver impacts the accuracy.
		\item An environment full of objects (like chair, table etc) reduces the accuracy due to multipath reflections.
		\item Cannot be used with crowded WiFi bandwidth. For example, 2.5 GHz band is crowded compared to 5 GHz band.
		\item User demographics are not mentioned.
	\end{itemize}
\end{frame}


\begin{frame}{Summary}
	\begin{itemize}
		\item WiFinger exploits the ubiquitous WiFi signals to sense f
		inger-grained gestures.
		\item The novelty of the system is the ability to extract fine-grained information from the CSI.
		\item WiFinger achieves an average recognition accuracy of 90.4\% per user.
	\end{itemize}

\end{frame}

%----------------------------------------------------------------------------------------
%	 CLOSING/SUPPLEMENTARY SLIDES
%----------------------------------------------------------------------------------------

\appendix

\begin{frame}{References}
	%\nocite{*} % Display all references regardless of if they were cited
	\bibliography{example.bib}
	\bibliographystyle{plain}
\end{frame}

%------------------------------------------------

\begin{frame}{Backup Slide}
	
\end{frame}

%----------------------------------------------------------------------------------------

\end{document}
